\textbf{乘法阶}

\textbf{定义}
在
\[
\gcd(a,m)=1
\]
时，$a$ 在模 $m$ 下的乘法阶记为 $\operatorname{ord}_m(a)$，定义为最小正整数 $k$，使得
\[
a^k\equiv 1\pmod m.
\]

\textbf{基本性质}
\begin{itemize}
\item $\operatorname{ord}_m(a)$ 一定存在。
\item $\operatorname{ord}_m(a)\mid \varphi(m)$。
\item $a$ 是模 $m$ 的原根，当且仅当 $\operatorname{ord}_m(a)=\varphi(m)$。
\end{itemize}

\textbf{常用求法}
\begin{itemize}
\item 先求 $\phi=\varphi(m)$。
\item 分解 $\phi$ 的不同质因子。
\item 从 $\phi$ 开始，对每个质因子 $p$ 尝试把当前答案除以 $p$。
\item 若
\[
a^{ans/p}\equiv 1\pmod m
\]
则保留这个缩小。
\end{itemize}

\textbf{快记}
乘法阶 = “$\varphi(m)$ 的最小可行因子”。
